% ============================================================
% Formulario, argomento 2: calcolo combinatorio (dispensa, capitolo 2)
% ============================================================
\argomento{Calcolo combinatorio}{Dispensa: capitolo 2}

\voce{fondamentale}{Principio fondamentale}{$k$ passi, $m_i$ scelte al passo $i$}
  {N = m_1 \cdot m_2 \cdots m_k = \prod_{i=1}^{k} m_i}
  {\variante{prodotto cartesiano}{\card{A_1 \times A_2 \times \cdots \times A_k} = \prod_{i=1}^{k} \card{A_i}}
   \variante{$k$ volte lo stesso insieme}{\card{A^k} = \card{A}^k}}

\voce{disprip}{Disposizioni con ripetizione}{ordine: sì $\cdot$ ripetizioni: sì}
  {K_{o,r} = n^k}
  {\variante{stringhe binarie lunghe $k$}{2^k}}

\voce{dispsemplici}{Disposizioni semplici}{ordine: sì $\cdot$ ripetizioni: no}
  {K_{o,nr} = n\,(n-1) \cdots (n-k+1) = \prod_{i=0}^{k-1} (n - i) = \frac{n!}{(n-k)!} \qquad (k \leq n)}
  {\variante{se $k > n$}{K_{o,nr} = 0}}

\voce{permutazioni}{Permutazioni}{tutti gli $n$ oggetti in fila}
  {n! = n\,(n-1) \cdots 2 \cdot 1}
  {\variante{è il caso $k = n$ delle disposizioni semplici}{K_{o,nr} = \frac{n!}{0!} = n!}
   \variante{con oggetti uguali ripetuti $n_1, \ldots, n_r$ volte \oltre}{\frac{n!}{n_1!\, n_2! \cdots n_r!}}}

\voce{fattoriale}{Fattoriale}{}
  {n! = \prod_{i=1}^{n} i = 1 \cdot 2 \cdots n \qquad 0! = 1}
  {\variante{ricorsione}{n! = n \cdot (n-1)!}
   \variante{rapporto ($m \leq n$)}{\frac{n!}{m!} = n\,(n-1) \cdots (m+1)}}

\voce{combinazioni}{Combinazioni}{ordine: no $\cdot$ ripetizioni: no\\ sottoinsiemi di $k$ elementi}
  {C_{n,k} = \binom{n}{k} = \frac{n!}{k!\,(n-k)!} = \frac{n\,(n-1) \cdots (n-k+1)}{k!}}
  {\variante{da due gruppi: $k_1$ fra $n_1$ e $k_2$ fra $n_2$}{\binom{n_1}{k_1} \binom{n_2}{k_2}}}

\voce{binomiale}{Proprietà del coefficiente binomiale}{}
  {\binom{n}{0} = \binom{n}{n} = 1 \qquad \binom{n}{1} = n \qquad \binom{n}{k} = \binom{n}{n-k}}
  {\variante{regola di Stifel ($1 \leq k \leq n-1$)}{\binom{n}{k} = \binom{n-1}{k-1} + \binom{n-1}{k}}}

\voce{combrip}{Combinazioni con ripetizione}{ordine: no $\cdot$ ripetizioni: sì\\ (nel corso non servono)}
  {\binom{n+k-1}{k}}
  {}

\voce{stringhe}{Stringhe con $k$ uni}{sequenze di $n$ prove con $k$ successi}
  {\frac{n!}{k!\,(n-k)!} = \binom{n}{k}}
  {}

\voce{newton}{Binomio di Newton}{}
  {(a + b)^n = \sum_{k=0}^{n} \binom{n}{k}\, a^k\, b^{n-k}}
  {\variante{$n = 2$ e $n = 3$}{(a + b)^2 = a^2 + 2ab + b^2, \quad (a + b)^3 = a^3 + 3a^2 b + 3ab^2 + b^3}}

\voce{sottoinsiemi}{Numero dei sottoinsiemi}{insieme delle parti di $A$, $\card{A} = n$}
  {\card{\mathcal{P}(A)} = \sum_{k=0}^{n} \binom{n}{k} = 2^n}
  {}
